\begin{abstract}
Bài báo này đề xuất phương pháp mạng từ trở sinh lưới trong không gian ba chiều (3D-MBGRN) dựa trên hệ tọa độ trụ phục vụ phân tích trường điện từ phi tuyến của các cấu trúc máy điện quay đối xứng, với mô hình kiểm chứng là động cơ nam châm vĩnh cửu hướng trục (AFPM). Phương pháp thực hiện rời rạc hóa miền tính toán dựa trên cấu trúc lưới lục diện đồng nhất kết hợp kỹ thuật dịch chỉ số (index-shifting), cho phép mô tả chuyển động quay của rotor mà không cần tái tạo lưới hay nội suy ranh giới. Bộ giải lặp phi tuyến được tích hợp thuật toán suy giảm thích nghi hệ số thư giãn vật liệu (Adaptive Material Decay) kết hợp cơ chế quay lui nghiệm (backtracking), giảm thiểu rủi ro nghiệm số học rơi vào bẫy trì trệ hội tụ. Kết quả kiểm chứng được đối chiếu với phương pháp phần tử hữu hạn ba chiều (3D-FEM) ghi nhận sự tương đồng về dạng sóng transient và phổ hài của các đại lượng điện từ cốt lõi ở cả trạng thái không tải và có tải định mức, với độ lệch giữa hai phương pháp dưới $3.1\%$. Khảo sát hội tụ lưới NRMSE cho thấy các nghiệm tìm được của phương pháp đề xuất hội tụ ổn định theo mật độ ô lưới. Kích thước ma trận hệ thống (DoF) của 3D-MBGRN cũng nhỏ hơn khoảng $1.4~\text{lần}$ so với 3D-FEM ở quy mô số lượng phần tử tương đương, giúp giảm tiêu tốn bộ nhớ RAM và rút ngắn thời gian tính toán transient nhanh hơn từ $6$ đến $16~\text{lần}$ . Các kết quả đạt được thể hiện khả năng ứng dụng của phương pháp 3D-MBGRN như một công cụ tính toán bổ trợ hoặc thay thế cho 3D-FEM trong các quy trình thiết kế và tối ưu hóa máy điện đa tham số.

\textbf{Keywords---} 3D Mesh-Based Reluctance Network (3D-MBGRN), Finite Element Method (FEM), Axial Flux Permanent Magnet Machine (AFPM), Periodic Boundary Conditions, Hexahedral Elements, Numerical Convergence, Material Relaxation.
\end{abstract}   